\section{Claim \#5}

\paragraph{Claim Statement} ``Generative AI raises knowledge-worker productivity by 25--40\% on tasks within its capability frontier, with little benefit outside it'' \citep{dellacqua2023jagged}.

\paragraph{Construct Identification} The claim's object is criterion-shaped. ``Knowledge-worker productivity'' is operationalized as task-completion rate, speed, and expert-rated quality on a fixed task suite drawn from industrial-engineering measurement traditions. The AI Construct Lexis \citep{aiconstructlexis} treats productivity as a downstream-utility construct whose validity hinges on the population and task domain in which it is measured. The headline reach to ``knowledge workers'' generally, rather than to BCG consultants on consulting-style tasks, is the inferential leap that has to bear the most weight.

\subsection{Content Validity}
\begin{itemize}
    \item \textbf{Rating:} Partial
    \item \textbf{Confidence:} 4
    \item \textbf{Written Argument:} The 18-task experimental suite was constructed by BCG and the authors to span inside- and outside-frontier work, with strong content validity for \emph{consulting} knowledge work such as memos, slides, ideation, and problem structuring. Coverage of other knowledge-work domains such as software engineering, clinical reasoning, legal analysis, and scientific research is zero by design. There is a systemic gap that the tasks selected exercise document- and ideation-heavy cognition where LLMs are strong, while embodied, empirical, and procedural work is absent from the sample \citep{salaudeen2025measurement}.
\end{itemize}

\subsection{Construct Validity}
\begin{itemize}
    \item \textbf{Rating:} Strong
    \item \textbf{Confidence:} 5
    \item \textbf{Written Argument:} Productivity is operationalized cleanly via three standard sub-measures (completion, speed, expert-rated quality) with double-blinded human rating. \emph{Structural validity} is met because the construct is industrial-engineering-standard and the three sub-measures cover its expected dimensions. \emph{Convergent validity} is met because the three measures move together across conditions in the study. \emph{Discriminant validity} is met because the inside-versus-outside frontier inversion (gain on inside tasks, loss on outside tasks) provides rare and credible evidence that the ``frontier'' sub-construct is identifiable rather than asserted after the fact. Construct validity within the experimental setting is unusually well established for this literature \citep{salaudeen2025measurement}.
\end{itemize}

\subsection{Criterion / Predictive Validity}
\begin{itemize}
    \item \textbf{Rating:} Partial
    \item \textbf{Confidence:} 3
    \item \textbf{Written Argument:} The study is a within-subjects field experiment with real consultants on real-shaped tasks, providing strong immediate concurrent evidence: randomized assignment, expert-rated quality, and the inside/outside-frontier inversion ($-19$pp on outside-frontier tasks). Predictive validity over weeks or months is not measured, and that includes novelty effects, deskilling, and over-reliance dynamics. The 25--40\% headline is a point-in-time effect. Longer-horizon dynamics are unstudied, and the criterion evidence is strong only on the immediate timescale \citep{salaudeen2025measurement}.
\end{itemize}

\subsection{Consequential Validity}
\begin{itemize}
    \item \textbf{Rating:} Partial
    \item \textbf{Confidence:} 4
    \item \textbf{Written Argument:} The study has been cited extensively in workforce-policy and hiring discussions, and the frontier nuance often gets dropped in citation chains, leaving a bare ``AI raises productivity 25--40\%'' that drives premature labor decisions. Consequential validity is partial. The original paper hedges appropriately, but the headline form of the claim invites loss of nuance, and the gap between the within-frontier finding and over-frontier deployment is the canonical consequential-validity risk \citep{salaudeen2025measurement}. The study itself is responsible. Its public-discourse footprint is less so.
\end{itemize}

\subsection{External validity}
\begin{itemize}
    \item \textbf{Rating:} Weak
    \item \textbf{Confidence:} 5
    \item \textbf{Written Argument:} \textbf{[MOST RELEVANT DIMENSION]} External validity is the linchpin because construct and criterion are internally well established. The entire question is whether the consulting-on-consulting-tasks finding generalizes to ``knowledge workers'' broadly. The study uses a single-firm sample (BCG), a single occupation (management consultant), a single tool (GPT-4 vintage), and a fixed task suite. The mechanism the headline invites the reader to extrapolate, namely a domain-general productivity uplift, is supported neither by sample diversity nor by independent replication in adjacent occupations. The qualitative \emph{frontier pattern} (intervention helps inside, hurts outside) is more transferable than the magnitude (25--40\%), because magnitudes depend on task-distribution base rates while the qualitative pattern reflects general LLM capability boundaries. The claim conflates a domain-specific magnitude with a domain-general pattern, and per Salaudeen et al., this is the core external-validity threat \citep{salaudeen2025measurement}.
\end{itemize}

\subsection{Overall Validity Judgment}
\begin{itemize}
    \item \textbf{Rating:} Partially Valid
    \item \textbf{Confidence:} 4
    \item \textbf{Written Argument:} The internal criterion evidence is strong, featuring a randomized field experiment and expert-rated quality, and the construct of productivity is cleanly operationalized. However, the critical failure point is external validity. The 25--40\% magnitudes apply specifically to BCG consultants executing consulting tasks, not to knowledge workers generically. Extrapolating these specific gains to entirely different domains constitutes a severe consequential validity risk, inviting premature workforce restructuring based on non-transferable metrics. The core generalizable finding is the boundary of the frontier itself, not the exact performance magnitude. \\
    \emph{Defensible reformulation}: ``In a randomized field experiment with management consultants at BCG on a fixed suite of tasks, generative AI use increased productivity on tasks inside the AI capability frontier (12--40\% gains) and decreased quality outside it ($-19$pp). Generalization of these magnitudes to other domains requires separate empirical work, while the qualitative frontier pattern is a more transferable finding.''
    \item \textbf{Peer Prediction (BTS):}

    \begin{tabular}{ll}
        Valid as stated: & 25\% \\
        Partially valid: & 60\% \\
        Not valid: & 15\% \\
        \end{tabular}\\[4pt]
\end{itemize}
